

\documentclass[]{kyvernitis-resume}
\fullname{Manish Dhakal}
\jobtitle{Computer Vision . Large Language Models }

\newcommand{\lbf}[1]{
    \contour{black}{\emph{#1}}
}

\usepackage[outline]{contour}% http://ctan.org/pkg/contour

\begin{document}
\resumeheader
% {\phone{+977 986-0687860}}
{\website{manishdhakal.github.io}}
{\linkedin{manishdhakal521}}
{\email{mdhakal3@gsu.edu}}
{\github{manishdhakal}}
{\gscholar{lhGdC2IAAAAJ}}
{\location{Atlanta, GA, USA}}
{\medium{manishdhakal}}

\begin{section}{Research}
Applied deep learning for computer vision and LLMs with the motive of \lbf{model compression and efficient learning.}
Publications at premier venues, such as \lbf{WACV, MICCAI, ACCV, IJCNLP-AACL,} etc.
\end{section}

\begin{section}{Education}
    \begin{subsection}{Ph.D. in Computer Science}{Georgia State University}{Graduate School}{August 2024 -- 2029*}
        \item[] {\textit{Supervisor:} \lbf{\href{https://ding1.com/}{Yi Ding, Ph.D.}}}
        \item \lbf{M.S. Degree} to be awarded in the summer of 2026.
        \item Deep learning projects on \lbf{2D/3D computer vision and large language models (LLMs),} funded by the US Department of Defense (DoD).
        \item Research in \lbf{efficient machine learning} and \lbf{model compression.}
        \italicitem{Grade: 3.9}
    \end{subsection}

    
    \begin{subsection}{Bachelor in Computer Engineering}{Pulchowk Campus, Institute of Engineering, Tribhuvan University}{Undergraduate School}{November 2017 -- April 2022}
        \item[] {\textit{Thesis Supervisor:} \lbf{\href{https://scholar.google.com/citations?user=ZB6i0OKaPo4C}{Prof. Subarna Shakya}}}
        \item \lbf{Ranked 11\textsuperscript{th} (top 0.1\%)} in the entrance exam, competing with \lbf{15,000+ candidates}, received a full scholarship for undergraduate studies.
        \item Gained knowledge about \lbf{significant CS courses} like AI, Image Processing, Data Structure \& Algorithm, DBMS, Software Engineering, and so on.
        \item \textit{Thesis:} Automatic speech recognition for \lbf{low-resourced Nepali language} which was later presented at an IEEE conference.
            % \italicitem{Thesis: The Art of Last-Minute Panic and Its Impact on Productivity}
            % \italicitem{Successfully procrastinated and panicked at the last moment}
            % \italicitem{Final Grade: 420/420}
    \end{subsection}
    
\end{section}

\begin{section}{Ongoing Projects}
    \begin{subsection}{Knowledge Distillation}{LLMs}{}{}
        \item Developing a robust knowledge distillation method utilizing a \lbf{teacher-student paradigm} to transfer complex knowledge efficiently.

        \item Engineering a feature-transfer mechanism between the teacher and student models' intermediate layers to capture and leverage rich, internal representations, ensuring robustness and high fidelity in model compression and knowledge transfer.
    \end{subsection}

    \begin{subsection}{LLM-Guided Adaptation for Vision-Language Models}{LLMs, VLMs}{}{}
    \item \lbf{Substituting visual prompts} with textual descriptions to enable pure LLMs to handle vision-language tasks, thereby mitigating the need for standard Vision-Language Models (VLMs).
    \item LLM-centric VLM adaptation method where the LLM itself provides preference feedback to optimize the textual descriptions.
    \end{subsection}
    
    \begin{subsection}{Parameter-Efficient Fine-tuning}{Point Cloud}{}{}
        \item \lbf{Graph Features Tuning (GFT)}, a Parameter-Efficient Fine-Tuning (PEFT) method specifically for point cloud data, reducing computational and memory costs compared to earlier approaches (\textit{accepted at WACV}).
        \item Achieved competitive performance on object classification and segmentation tasks while substantially \lbf{reducing the number of trainable parameters} by integrating graph features into deeper layers via skip connections and efficient cross-attention.
    \end{subsection}


\end{section}


\begin{section}{Publications}
    \begin{subsection}
    {2022-Present}{}{}{\href{https://scholar.google.com/citations?user=lhGdC2IAAAAJ}{\underline{Scholar}} Citations: 100+ \& h-index: 4}
        \item \lbf{\underline{Dhakal, M.}}, Dasari, V. R., Sunderraman, R., \& Ding, Y. (2026, March). GFT: Graph feature tuning for efficient point cloud analysis. In \textit{WACV.} [\href{https://arxiv.org/abs/2511.10799}{Paper}, \href{https://github.com/manishdhakal/GFT}{Code}]
        
        \item Mim, S. T., Morris, J., \lbf{\underline{Dhakal, M.}}, Xiu, Y., Gorlatova, M., \& Ding, Y. (2025, December). Can a unimodal language agent provide preferences to tune a multimodal vision-language model? In \textit{IJCNLP-AACL Findings.}
        
        \item Budathoki, A., \& \lbf{\underline{Dhakal, M.}} (2025). Adversarial robustness analysis of vision-language models in medical image segmentation. \textit{arXiv preprint arXiv:2505.02971.} [\href{https://arxiv.org/abs/2505.02971}{Paper}, \href{https://github.com/anjilab/secure-private-ai}{Code}]
        
        \item Adhikari, R., Thapaliya, S., \lbf{\underline{Dhakal, M.}}, \& Khanal, B. (2024). TuneVLSeg: Prompt tuning benchmark for vision-language segmentation models. In \textit{ACCV (Oral).} [\href{https://openaccess.thecvf.com/content/ACCV2024/html/Adhikari_TuneVLSeg_Prompt_Tuning_Benchmark_for_Vision-Language_Segmentation_Models_ACCV_2024_paper.html}{Paper}, \href{https://github.com/naamiinepal/tunevlseg}{Code}]
        
        \item \lbf{\underline{Dhakal, M.}}, Adhikari, R., Thapaliya, S., \& Khanal, B. (2024). Finetuning vision-language segmentation efficiently with lightweight blocks. In \textit{MICCAI.} [\href{https://arxiv.org/abs/2405.06196}{Paper}, \href{https://github.com/naamiinepal/vlsm-adapter}{Code}]
        
        \item Adhikari, R.$^{*}$, \lbf{\underline{Dhakal, M.$^{*}$}}, Thapaliya, S.$^{*}$, Poudel, K., Bhandari, P., \& Khanal, B. (2023, October). Synthetic Boost: Leveraging synthetic data for enhanced vision-language segmentation in echocardiography. In \textit{ASMUS @ MICCAI.} [\href{https://arxiv.org/abs/2309.12829}{Paper}, \href{https://github.com/naamiinepal/synthetic-boost}{Code}]
        
        \item Poudel, K.$^{*}$, \lbf{\underline{Dhakal, M.$^{*}$}}, Bhandari, P.$^{*}$, Adhikari, R.$^{*}$, Thapaliya, S.$^{*}$, \& Khanal, B. (2024). Exploring transfer learning in medical image segmentation using vision-language models. In \textit{MIDL (Oral).} [\href{https://proceedings.mlr.press/v250/poudel24a.html}{Paper}, \href{https://github.com/naamiinepal/medvlsm}{Code}]
        
        \item \lbf{\underline{Dhakal, M.}}, Chhetri, A., Gupta, A. K., Lamichhane, P., Pandey, S., \& Shakya, S. (2022, July). Automatic speech recognition for the Nepali language using CNN, bidirectional LSTM, and ResNet. In \textit{Proceedings of the International Conference on Inventive Computation Technologies (ICICT).} IEEE. [\href{https://arxiv.org/abs/2406.17825}{Paper}, \href{https://github.com/manishdhakal/ASR-Nepali-using-CNN-BiLSTM-ResNet}{Code}]
        
    \end{subsection}
\end{section}

\begin{section}{Work Experience} 
    \begin{subsection}{Nepal Applied Mathematics and Informatics Institute for research (NAAMII)}{Research Assistant}{Lalitpur, Nepal }{April 2022 -- June 2024}
        \item[] {\textit{Supervisor:} \lbf{\href{https://scholar.google.com/citations?user=ZfaUCG5h3xsC}{Bishesh Khanal, Ph.D.}}}
        \item Developed skills for \lbf{object detection and segmentation} tasks on 2D medical images and explored their multi-modal approach ( esp. \lbf{vision-language models} ); also worked on segmentation with \lbf{3D mesh} data.
        \item Demonstrated \lbf{strong skills in writing scientific manuscripts,} with multiple papers submitted for review, showcasing the ability to \lbf{communicate methodologies, results, and implications} effectively.
        \item Ensured \lbf{reproducibility and modularity in ML projects} by implementing robust methodologies and practices, allowing for the transparent and replicable programming of the projects.
    \end{subsection}
\end{section}

\begin{section}{Past Projects}
\begin{subsection}{Lower Limb Segmentation}{Medical Imaging}{July 2023 -- September 2023}{}
    \item[] \textit{Supervisor:} \lbf{\href{https://scholar.google.com/citations?user=RuVA5rUAAAAJ}{Taman Upadhaya, Ph.D.}}
    \item Conducted \lbf{training experiments} of different deep learning models on the remote server to segment three bones -- knee, pelvis, and ankle -- from CT scans of the lower limbs of patients.
    \item \lbf{Deployed a robust Python rest API} on the remote server for the segmentation request from a client, with a pipeline including pre-processing, inference, and post-processing steps.
    \item Ensured \lbf{interoperability, reproducibility, and understandability} of the deployed application using Docker, and well-structured documentation and comments.
\end{subsection}

\begin{subsection}{Vision Language Segmentation Models (VLSMs) for Medical Images}{Medical Imaging}{May 2023 -- August 2024}{}
    \item Reported zero-shot and finetuned segmentation performance of \lbf{4 VLSMs} on \lbf{11 medical datasets} using \lbf{9 types of prompts} derived from \lbf{14 attributes}, prompts are given as text conditioning information.
    \item Worked with \lbf{encoder-decoder architecture} to generate binary segmentation masks for VLSMs.
    \item Tested the compatibility of the VLSMs (such as \lbf{CLIPSeg and CRIS}) pre-trained for open-domain images with medical images.
\end{subsection}

\begin{subsection}{Object Detection in 2D Orthopantomogram (OPG) Images}{Dental Imaging}{September 2022 -- August 2024}{}
    
    \item Critically analyzed the \lbf{literature and state-of-the-art} models for different segmentation and detection tasks on radiology images of dentistry and their inadequacy.
    \item Designed and developed the \lbf{data annotation tool} for object detection over 2D OPG images.
    \item Working on identification and localization of dental \lbf{anatomical structures and abnormalities} while benchmarking with existing methods like \lbf{YOLO, RetinaNet, RCNN, and FastRCNN}.
\end{subsection}


\begin{subsection}{Segmentation in 3D Teeth Scan}{MICCAI Challenge 2022}{Summer 2022}{}
    \item Learned about the representation and preprocessing of \lbf{3D mesh and point cloud} data. 
    \item Benchmarked with different 3D point cloud segmentation models such as \lbf{Pointnet/++ and DeltaConv}.
\end{subsection}

\begin{subsection}{Nepali AutoComplete and LM}{\href{https://github.com/gaudelbijay/NepaliAutoCompleteAndLM}{Open Source Project}}{August 2020 -- October 2020}{}
    \item Designed and trained \lbf{language model of Nepali (ie. Devnagari transcript)} for the text auto-complete system.
    \item Programmed the \lbf{pre-processing pipeline} to remove the non-Nepali characters from the dataset.
\end{subsection}

\begin{subsection}{Super-Resolution with GAN (SRGAN)}{\href{https://github.com/manishdhakal/SuperResolution}{Open Source Project}}{May 2020 -- August 2020}{}
    \item Implemented \lbf{open source model} of \lbf{SRGAN} with Keras/TensorFlow.
    \item Developed the \lbf{understanding of generator and discriminator} in GAN-based generative models.
\end{subsection}

\end{section}


\begin{section}{Teaching}
    \begin{subsection}{Community Eye, ENT \& Rehabilitation Center (CEERS)}{Trainer}{Bhaktapur, Nepal}{June 2023 -- June 2024}
        \item \lbf{Training a group of interns} to develop medical imaging applications with the use of ML.
        \item Instructing and guiding them about ML through activities like \lbf{paper reading sessions, lecture-lab sessions,} and \lbf{topic presentations}. 
    \end{subsection}
    
    \begin{subsection}{4th Annual Nepal AI School (ANAIS)}{Teaching Assistant}{Kathmandu, Nepal}{May 2023 -- June 2023}
        \item \lbf{Guided participants} through a series of labs related to \lbf{neural networks, transformers, federated learning,} \lbf{graph neural networks, active learning,} and so on.
        \item \lbf{Mentored three groups} during the 10-day \lbf{machine learning hackathon} (namely, Hack-a-Dev).
    \end{subsection}

    \begin{subsection}{Software Fellowship, Locus 2021}{Programming Instructor}{Online}{Summer 2021}
        \item Provided tutoring on \lbf{software development life cycle} and assisted participants with \lbf{software documenta}-\lbf{tion} and \lbf{library/framework installation}.
        \item Taught participants about \lbf{API development for web applications,} emphasizing its concepts, best practices, and usage.
    \end{subsection}
\end{section}

    


\sectiontable{Achievements/Services}{
    \entry{2026}{\lbf{MICCAI Reviewer}}
    \entry{2025}{\lbf{AAAI Reviewer}}
    \entry{2024}{\lbf{Presidential Fellowship,} GSU (Only 6\% of the doctoral students @ GSU).}
    \entry{2024}{\lbf{LMIC Travel Grant} by the conference of MICCAI.}
    \entry{2017}{\lbf{Undergraduate} funded by the Government of Nepal.}
    % \entry{Winner}{LogPoint CTF, Cybersecurity Competition.}
    % \entry{Finalist}{Spirathon 2020, TechXperience Nepal.}
}

\sectiontable{Technical skills}{
    % \entry{Machine Learning}{Projects for detection and segmentation task while maintaining \lbf{reproducibility} \lbf{and modularity}; integrating \lbf{open source models} for benchmarking. Proficiency in using libraries and frameworks like \lbf{NumPy, Pandas, PyTorch, and TensorFlow}.}\\
    \entry{Deep Learning}{\lbf{Unimodal and multimodal learning}, large language models (LLMs), model quantization, efficient learning, reinforcement learning fine-tuning, 2D/3D vision, convolution, and transformers.}\\
    
    \entry{Writing}{\lbf{Knowledge synthesis} from the existing literature, \lbf{writing scientific documents} and \lbf{manuscripts} with \lbf{LaTex}, and \lbf{communicating the results} to the community with transparency.}\\
    
    % \entry{Web Development}{Competence in creating well-documented backend applications with relational databases using frameworks like \lbf{Django, FastAPI, and NodeJS}. Adept at client-side programming with \lbf{ReactJS.}}\\
    
    \entry{Remote Server}{Able to work with \lbf{remote Linux machines} for coding and project deployment using \lbf{SSH, shell script, tmux, Nginx, and Docker}.}\\
    
}

% \begin{section}{Extracurricular Activities}
%     \begin{subsection}{DataRush (AI and Data Science Competition, Locus)}{Co-ordinator}{Spring 2022}{}
%         \item \lbf{Call for sponsors}, maintained communication and coordination between sponsors, participants, mentors, and other organizers.
%         \item Made the \lbf{budget planning}, prepared the \lbf{event's rule book}, planned the \lbf{event structure}, and ensure the smooth operations of the event.
%         \item \lbf{Tested and validated} the machine learning models and their \lbf{solutions} submitted by the participants.
%     \end{subsection}
% \end{section}



% \sectiontable{Certifications}{
%     \entry{Stanford University}{\href{https://www.coursera.org/account/accomplishments/verify/5ZN4D4J4VA3E}{Machine Learning}}
%     \entry{DeepLearning.AI}{\href{https://www.coursera.org/account/accomplishments/specialization/NHRM9UQJYAFJ}{Deep Learning Specialization}}
%     \entry{DeepLearning.AI}{\href{https://www.coursera.org/account/accomplishments/specialization/LWTPDKJWXK84}{AI for Medicine}}
%     \entry{DeepLearning.AI}{\href{https://www.coursera.org/account/accomplishments/specialization/ZCC2D7N3WCSC}{Natural Language Processing}}
%     \entry{DeepLearning.AI}{\href{https://www.coursera.org/account/accomplishments/specialization/H93N38K9FQ8V}{TensorFlow in Practice}}
%     \entry{University of Alberta}{\href{https://www.coursera.org/account/accomplishments/specialization/5D5T9HQCNVEY}{Reinforcement Learning}}
% }


\begin{section}{References (Research Advisors)}
    \begin{subsectionnobullet}
        {Yi Ding, Ph.D.}{ {\normalfont Assistant Professor,} Department of Computer Science and Software Engineering, Auburn University}{}{}
        \item{\href{mailto:yiding@gsu.edu}{\ttfamily yiding@gsu.edu}}
    \end{subsectionnobullet}
    \begin{subsectionnobullet}
        {Raj Sunderraman, Ph.D.}{ {\normalfont Professor,} Department of Computer Science, Georgia State University (GSU)}{}{}
        \item{\href{mailto:rsunderraman@gsu.edu}{\ttfamily rsunderraman@gsu.edu}}
    \end{subsectionnobullet}
    \begin{subsectionnobullet}
        {Bishesh Khanal, Ph.D.}{ {\normalfont Research Director,} Nepal Applied Mathematics and Informatics Institute for research (NAAMII)}{}{}
        \item{\href{mailto:bishesh.khanal@naamii.org.np}{ \ttfamily bishesh.khanal@naamii.org.np}}
    \end{subsectionnobullet}

\end{section}

% \sectiontable{Soft skills}{
%     \entry{Procrastination}{Procrastination at an expert level}
%     \entry{Avoiding Responsibility}{Outstanding ability to avoid meetings and responsibilities}
%     \entry{Sarcasm}{Fluent in sarcasm and irony}
% }

% \sectiontable{Awards}{
%     \entry{World's Best Napper}{International Association of Snoozers \hfill \textit{2022}}
%     \entry{Most Creative Excuses}{Academy of Procrastinators \hfill \textit{2021}}
%     \entry{Gold Medal in Avoiding Responsibilities}{Olympics of Slackers \hfill \textit{2020}}
% }

\end{document}
